% CounterPrune draft. Generic article class so it compiles anywhere; swap for the venue class (cvpr/neurips/iclr).
% Every number comes from docs/results-and-execution-plan.md, which is recomputed from results/ by tools/analysis/.
% Sections marked TODO depend on the LLaVA-NeXT follow-up run (HANDOFF §7 item 3).
\documentclass[10pt]{article}
\usepackage[margin=1in]{geometry}
\usepackage{amsmath,amssymb,booktabs,multirow,graphicx,xcolor,hyperref,natbib}
\newcommand{\method}{CounterPrune}
\newcommand{\todo}[1]{\textcolor{red}{[TODO: #1]}}
\newcommand{\sig}{$^{*}$}

\title{\method: Training-Free Per-Image Visual-Token Exit for Vision-Language Models}
\author{Anonymous}
\date{}

\begin{document}
\maketitle

\begin{abstract}
Visual tokens dominate the prefill cost of LLaVA-style vision-language models (VLMs): a single image contributes 576
tokens in LLaVA-1.5 and up to 2{,}928 in LLaVA-NeXT, more than 95\% of the sequence. Attention-guided token pruning
reduces this cost but must materialise attention weights, which breaks fused attention kernels, and prior analyses show
that the ranking criterion matters less than how many layers see the full token set. We propose \method, a
training-free method that removes \emph{all} visual tokens at a layer chosen \emph{per image} by a cheap signal read
during the forward pass: either the logit-lens stability of the last token (KL divergence between consecutive layers'
next-token distributions) or a one-row counterfactual probe (the change in the last token's attention update when the
visual keys are removed, computed from the layer's KV cache). Neither signal needs attention matrices, so \method{}
runs on SDPA/FlashAttention with a standard KV cache; later layers simply cache fewer tokens. With one global threshold
chosen by a held-out calibration protocol, \method{} keeps the unpruned model's answers with no statistically
significant accuracy change on POPE, MME, GQA and TextVQA at 68\% of baseline prefill FLOPs on LLaVA-1.5-7B, whereas
the dataset-level fixed exit layer (VTW) selected by the same protocol needs 88\%. The exit layer tracks a per-example
oracle (Spearman 0.5--0.6 on GQA and TextVQA) and separates question types. Measured prefill latency follows FLOPs, the
probe costs under 20\,ms, and \method{} composes with FastV-style token pruning on the same engine. The same threshold
transfers to LLaVA-NeXT-7B on POPE, MME and GQA; OCR-heavy TextVQA needs the higher fidelity target. \todo{update
after the NeXT follow-up run}
\end{abstract}

\section{Introduction}
\todo{write; key points below}
\begin{itemize}
\item Prefill cost is dominated by visual tokens; attention is quadratic in length.
\item Attention-guided pruning (FastV, SparseVLM, PDrop) ranks tokens with decoder attention: needs eager attention,
  and the ranking criterion matters little (FEATHER, Wen et al.; our InfoScore analysis, \S\ref{sec:analysis}).
\item Visual-token exit (VTW) is kernel-compatible but dataset-level; per-instance exits (DyVTE, AdaVSkip) are trained.
\item Contributions: (i) a training-free per-image exit with two kernel-compatible signals; (ii) a calibration
  protocol that yields one global threshold and the fair VTW baseline; (iii) an SDPA + KV-cache engine in which exit,
  fixed exit and FastV are evaluated with measured latency; (iv) an analysis showing why attention-ranking quality
  does not matter (depth dominates) and that a per-image signal has to come from elsewhere.
\end{itemize}

\section{Related work}
\paragraph{Attention-guided token pruning.} FastV \citep{chen2024fastv} prunes low-attention visual tokens after a
fixed early layer; SparseVLM \citep{zhang2024sparsevlm} uses text-selected raters with progressive pruning; PyramidDrop
and STAR-Pro refine the schedule. All rank tokens with decoder attention, which requires materialising attention
weights and is incompatible with fused kernels.
\paragraph{Critiques.} \citet{endo2025feather} show that FastV's early-layer attention is dominated by RoPE-induced
positional bias and fails on localisation; \citet{wen2025token} find that random selection and pooling match elaborate
criteria on most benchmarks. Our layer-wise InfoScore analysis gives an information-theoretic account: attention in
layers 0--1 is nearly image-independent, yet ranking with the most content-adaptive layer yields no gain.
\paragraph{Visual-token exit.} VTW \citep{lin2024vtw} removes all visual tokens after a dataset-level layer chosen by a
KL criterion on a sample. DyVTE \citep{dyvte2024} and AdaVSkip \citep{adavskip2026} make the decision per instance with
trained hypernetworks or routers. HiDrop combines late injection, progressive pruning and exit in a trained fixed
schedule. \method{} is, to our knowledge, the first training-free per-instance exit that is also compatible with fused
attention kernels and KV caching. \todo{re-check 2025--26 literature before submission}

\section{Method}
\label{sec:method}
\paragraph{Setting.} A decoder with $L$ layers processes a sequence of $n_0$ tokens, of which a contiguous block of
$n_v$ are visual. Prefill FLOPs per layer at length $n$ follow the FastV formula $4nd^2 + 2n^2d + 2ndm$ ($d$ hidden,
$m$ FFN size); we report FLOPs summed over layers at the actual per-layer sequence length, relative to the unpruned
model.
\paragraph{Per-image exit.} Before (or after) each layer $E$, a scalar signal $s_E$ is read from the last token. Once
$s_E \le \tau$ for $m$ consecutive layers and $E \ge E_{\min}$, all visual tokens are dropped: hidden states are sliced,
original position ids are kept, and layers $E{:}L$ attend over text tokens only. The KV cache of later layers holds
fewer tokens, so decoding needs no special masking.
\paragraph{Signal 1: logit-lens stability (lens).} $s_E = \mathrm{KL}\big(\mathrm{lens}_E \,\|\, \mathrm{lens}_{E-1}\big)$
where $\mathrm{lens}_E = \mathrm{softmax}(W_U\,\mathrm{norm}(h^{(E)}_{\text{last}}))$. Cost: one LM-head product
per layer on one token. Default $\tau=0.05$, $m=1$, $E_{\min}=1$ (the signal never fires before layer $\approx$17 in
practice).
\paragraph{Signal 2: one-row counterfactual probe (cf).} After layer $E$, read its cached keys/values, recompute the
last token's query, and form one attention row with and without the visual keys:
$s_E = \|W_O(\mathrm{Attn}(q, K, V) - \mathrm{Attn}(q, K_{\neg v}, V_{\neg v}))\| / \|h_{\text{last}}\|$.
Cost: one attention row and one $W_O$ product. Default $\tau=0.05$, $m=2$, $E_{\min}=12$.
\paragraph{Kernel compatibility.} Neither signal needs attention matrices. The engine runs stock layers with
\texttt{attention\_mask=None} (causal SDPA), explicit position ids and a \texttt{DynamicCache}; it reproduces the stock
model's logits to $|\Delta| \le 0.008$ in fp16 and its greedy text.
\paragraph{Calibration protocol.} One global $\tau$ per signal is the cheapest grid value whose pooled agreement with
the unpruned model on a calibration split reaches a target $T$. Applying the same rule to a fixed exit layer $E$ is
exactly VTW's dataset-level selection, which gives the fair fixed-exit baseline.
\paragraph{Composition with token pruning.} FastV's criterion (last-token attention over visual tokens at layer $K-1$)
is the same one-row attention, so FastV runs on the same engine and can precede the exit.

\section{Experiments}
\label{sec:exp}
\paragraph{Setup.} LLaVA-1.5-7B and LLaVA-NeXT-7B (Vicuna), fp16, SDPA, KV cache, batch 1, 2$\times$T4. POPE (900;
300 per split), MME (660, $\le$50 per category; two questions per image), GQA testdev-balanced (600), TextVQA (500, VQA
soft accuracy, no OCR tokens). Greedy decoding, full-answer scoring. We report accuracy, agreement with the unpruned
model's answers, relative prefill FLOPs and measured prefill latency; paired accuracy differences use the exact McNemar
test. Standard errors: $\pm$1.1 (POPE), $\pm$1.7 (MME), $\pm$2.0 (GQA), $\pm$2.2 (TextVQA) points.

\subsection{Main result: one global threshold}
\begin{table}[t]
\centering\small
\caption{LLaVA-1.5-7B, one global setting. Accuracy / agreement with the unpruned model / relative prefill FLOPs (\%).
\sig: significantly worse than the unpruned model (exact McNemar, $p<0.05$). Fixed-E rows are the VTW-style exit.}
\label{tab:main}
\begin{tabular}{lccccc}
\toprule
Method & POPE & MME & GQA & TextVQA & Mean FLOPs \\
\midrule
Unpruned & 88.6 / 100 / 100 & 74.4 / 100 / 100 & 60.2 / 100 / 100 & 50.8 / 100 / 100 & 100 \\
Fixed E16 & 88.8 / 99.6 / 52 & 74.7 / 98.8 / 53 & 53.0\sig / 73.3 / 53 & 17.1\sig / 27.0 / 53 & 52 \\
Fixed E24 & 88.6 / 100 / 76 & 74.2 / 98.6 / 76 & 59.2 / 93.0 / 76 & 48.5\sig / 82.8 / 76 & 76 \\
Fixed E28 & 88.6 / 100 / 88 & 74.2 / 98.9 / 88 & 59.8 / 96.7 / 88 & 50.2 / 93.0 / 88 & 88 \\
FastV K2 R0.5 & 88.2 / 95.0 / 54 & 73.6 / 95.6 / 55 & 57.0\sig / 88.2 / 55 & 51.1 / 86.8 / 55 & 55 \\
FastV K2 R0.75 & 84.7\sig / 88.6 / 32 & 72.6 / 90.3 / 33 & 54.3\sig / 78.2 / 33 & 49.9 / 76.8 / 33 & 33 \\
\midrule
\method{} lens $\tau$=0.05 & 88.6 / 99.8 / 59 & 75.2 / 98.6 / 60 & 59.3 / 95.5 / 73 & 49.9 / 89.8 / 80 & \textbf{68} \\
\method{} cf $\tau$=0.05 & 88.4 / 98.8 / 50 & 73.8 / 95.2 / 50 & 59.8 / 96.8 / 76 & 50.7 / 97.4 / 93 & 67 \\
\method{} cf $\tau$=0.025 & 88.4 / 99.9 / 67 & 74.4 / 98.5 / 67 & 60.2 / 99.3 / 89 & 50.8 / 100 / 98 & 80 \\
FastV K2 R0.5 + lens $\tau$=0.05 & 88.3 / 95.1 / 34 & 73.9 / 95.3 / 35 & 56.8\sig / 85.7 / 42 & 50.1 / 80.8 / 45 & 39 \\
\bottomrule
\end{tabular}
\end{table}
Table~\ref{tab:main} (results doc \S8E/\S8F): with $\tau=0.05$, neither signal changes accuracy significantly on any
benchmark. The only fixed layer that is equally safe everywhere is E28 at 88\% FLOPs; fixed E24 (76\%) is
significantly worse on TextVQA and E16 keeps only 27\% of TextVQA answers. Mean FLOPs for lens exit: 68\% (59--60\% on
POPE/MME, 73\% GQA, 80\% TextVQA). Fixed exit layers fail on tasks that need the image late; the per-image exit adapts.

\subsection{Calibration protocol and the VTW baseline}
\begin{table}[t]
\centering\small
\caption{Threshold chosen on a 20\% calibration split of every benchmark (pooled agreement $\ge T$), evaluated on the
remaining 80\%; 50 random splits. Same rule over fixed $E$ = VTW. Evaluation FLOPs / agreement / mean accuracy change.}
\label{tab:calib}
\begin{tabular}{llccc}
\toprule
$T$ & Family & Chosen (seeds) & Eval FLOPs & Eval agree / acc $\Delta$ \\
\midrule
95\% & Fixed (VTW) & E28 $\times$48, E24 $\times$2 & 87.6 & 97.0 / $-$0.30 \\
95\% & lens & $\tau$=0.05 $\times$28, 0.1 $\times$12, 0.0245 $\times$9, 0.01 $\times$1 & \textbf{68.0} & 95.7 / $-$0.25 \\
95\% & cf & $\tau$=0.05 $\times$50 & 66.9 & 96.9 / $-$0.33 \\
97\% & Fixed (VTW) & E28 $\times$50 (target met in 24/50) & 88.1 & 97.2 / $-$0.27 \\
97\% & lens & $\tau$=0.01 $\times$30, 0.0245 $\times$15, 0.05 $\times$5 & 73.7 & 97.3 / $-$0.33 \\
97\% & cf & $\tau$=0.025 $\times$26, 0.05 $\times$24 & 73.9 & 98.1 / $-$0.19 \\
98\% & Fixed (VTW) & E28 $\times$50 (target met in 5/50) & 88.1 & 97.2 / $-$0.27 \\
98\% & lens & $\tau$=0.01 $\times$46, 0.0245 $\times$4 & 75.5 & 97.7 / $-$0.39 \\
98\% & cf & $\tau$=0.025 $\times$49, 0.05 $\times$1 & 80.1 & 99.4 / $-$0.05 \\
\bottomrule
\end{tabular}
\end{table}
Table~\ref{tab:calib} (results doc \S8G): the protocol recovers $\tau=0.05$ for both signals at $T=95\%$ and places
\method{} 14--20 FLOPs points below the VTW-style fixed layer at every target, at equal or higher evaluation agreement.
Significant losses on the evaluation split occur in at most 6\% of seeds, the nominal rate. Leave-one-benchmark-out:
calibrated on POPE/MME/GQA at $T=95\%$, fixed E24 ($-$2.3\sig) and lens $\tau=0.2$ ($-$3.3\sig) fail on TextVQA while
cf $\tau=0.05$ holds ($-$0.5); at $T=98\%$ every family transfers.

\subsection{Matched fidelity}
At the agreement each adaptive point reaches, a fixed layer needs 8--13 FLOPs points more on GQA (lens $\tau=0.05$:
95.5\% at 73\% vs.\ fixed 84\%) and 3--4 on TextVQA; cf $\tau\le0.05$ and lens $\tau\le0.0245$ reach 92--100\%
agreement on TextVQA, which no fixed $E\le28$ reaches below 97\% FLOPs. On POPE and MME a per-benchmark-tuned E16 is as
good as any adaptive point: visual dependency is homogeneous there and there is nothing to adapt to.

\subsection{Adaptivity}
The exit layer correlates with a per-example oracle (the earliest layer after which dropping all visual tokens keeps
the baseline answer): Spearman $+0.51$ (GQA) and $+0.57$ (TextVQA) for cf, $+0.42$ / $+0.22$ for lens; the exit is at or
after the oracle for 91--99\% of examples. GQA open-ended ``query'' questions exit at layer 25.6 on average, ``choose''
22.0, ``compare'' 19.3, ``verify'' 18.6, ``logical'' 18.3. MME categories all exit at 17.7--19.7, like POPE.

\subsection{Latency and probe cost}
Measured prefill time follows FLOPs: lens $\tau=0.05$ 336\,ms at 59\% FLOPs vs.\ fixed E20 343\,ms at 64\% (POPE,
LLaVA-1.5, baseline 498\,ms). The residual of each method's prefill time against a linear FLOPs$\to$time fit of the
unpruned and fixed-exit runs measures probe cost: lens $+5$ to $+11$\,ms, cache-based cf $+21$\,ms (a recompute-based
cf probe costs $+77$\,ms), FastV-in-engine $+22$ to $+25$\,ms. On LLaVA-NeXT (prefill $\approx$2\,s) all probes are
below 1.5\% of prefill.

\subsection{Comparison and composition with FastV on the same engine}
FastV K2 R0.5 keeps 88\% of answers at 55\% FLOPs and is significantly worse on GQA; \method{} at 67--68\% keeps
92--94\% with no significant loss. At FastV's operating point the exit alone is worse (cf $\tau=0.1$: 86.5\%,
$-$2.5 points). The two compose: FastV K2 R0.5 + lens $\tau=0.05$ reaches 39\% FLOPs at $-$1.2 points, against FastV K2
R0.75 at 33\% and $-$3.1 points (significant on POPE and GQA). Uniform spatial subsampling of visual tokens (50\% from
layer 2) is an ablation, not part of the method: $-$2.3 points mean, $-$6.5 on TextVQA.

\subsection{Second model: LLaVA-NeXT-7B}
\todo{replace with the follow-up run (600 GQA / 500 TextVQA, $\tau\in\{0.01,0.0245,0.05,0.1\}$, cf 0.025/0.05, E16--28,
FastV R0.5/R0.75, FastV+lens)}. Current 300-example runs with the LLaVA-1.5 thresholds (results doc \S8F.3): lens
$\tau=0.05$ keeps 100 / 99.0 / 96.0\% of answers on POPE / MME / GQA at 59 / 61 / 70\% FLOPs with no significant change
(prefill $-$36\% on POPE), but is significantly worse on TextVQA ($63.8$ vs.\ $68.0$, $p=0.013$), where NeXT's higher
resolution makes the image matter deeper; cf $\tau=0.05$ holds everywhere (98.7\% agreement on TextVQA at 94\% FLOPs).
This mirrors the leave-one-benchmark-out result on LLaVA-1.5: a threshold calibrated on early-saturating tasks is too
loose for OCR-heavy tasks unless the fidelity target is $\ge$97--98\% or the cf signal is used.

\section{Analysis: why attention-ranking quality does not matter}
\label{sec:analysis}
\todo{condense results doc \S3--\S6 (InfoScore layers 0--1 image-independent, 2--3 most content-adaptive; ranking
with any layer gives the same accuracy; depth dominates; per-image attention entropy has AUC $\approx$0.5 for
predicting pruning damage; oracle headroom at R=0.9: 29\% / 34\% FLOPs vs.\ 59--60\% for K=17).}

\section{Limitations}
Per-image exits break batching (batch 1, or pad to the latest exit in the batch). The lens signal never fires before
layer $\approx$17, so savings are bounded at about 45\% on tasks that saturate early (cf fires from layer 12). FLOPs
and latency are prefill-only; decoding benefits only through the smaller KV cache. Timings are SDPA on T4;
FlashAttention-2 numbers need an Ampere GPU. Localisation benchmarks (RefCOCO) are not evaluated.

\bibliographystyle{plainnat}
\begin{thebibliography}{9}
\bibitem[Chen et al.(2024)]{chen2024fastv} L. Chen et al. An image is worth 1/2 tokens after layer 2: Plug-and-play inference acceleration for large vision-language models. ECCV 2024.
\bibitem[Zhang et al.(2024)]{zhang2024sparsevlm} Y. Zhang et al. SparseVLM: Visual token sparsification for efficient vision-language model inference. arXiv:2410.04417.
\bibitem[Endo et al.(2025)]{endo2025feather} M. Endo et al. Feather the throttle: Revisiting visual token pruning for vision-language model acceleration. ICCV 2025.
\bibitem[Wen et al.(2025)]{wen2025token} Z. Wen et al. Token pruning in multimodal large language models: Are we solving the right problem? ACL Findings 2025.
\bibitem[Lin et al.(2024)]{lin2024vtw} Z. Lin et al. Boosting multimodal large language models with visual tokens withdrawal for rapid inference. arXiv:2405.05803.
\bibitem[DyVTE(2024)]{dyvte2024} Accelerating multimodal large language models via dynamic visual-token exit and the empirical findings. arXiv:2411.19628.
\bibitem[AdaVSkip(2026)]{adavskip2026} Adaptive visual token skipping across layers for efficient MLLMs inference. arXiv:2609.15131.
\end{thebibliography}
\end{document}
